\documentclass[11pt]{article}
\usepackage[colorlinks=true,linkcolor=blue,urlcolor=blue]{hyperref}
\font\cap=cmcsc10
\setlength{\topmargin}{-0.15in}
\setlength{\oddsidemargin}{-0.15in}
\setlength{\evensidemargin}{-0.15in}
\setlength{\headheight}{0.1in}
\setlength{\headsep}{0.25in}
\setlength{\topskip}{0.1in}
\setlength{\textwidth}{6.5in}
\setlength{\textheight}{9.25in}
\setlength{\parindent}{0.0in}
\setlength{\parskip}{1em}
\usepackage{macros_cavag}
\usepackage{graphicx}
\usepackage{hyperref}
\usepackage{epsfig}
\pagestyle{empty}

\begin{document}

\today

Prof. Patricia Burchat\\
Stanford University\\
382 via Pueblo Mall\\
Varian 158\\
Stanford, CA 94305-4060\\

Dear Prof. Burchat:

Please accept the attached application for the postdoctoral
position in galaxy cluster studies advertised in the September 2007 edition
of the AAS Job Register. For my thesis, entitled {\textit{`Feedback,
Evolution, and Dynamics in Galaxy Clusters'}}, I am studying
the connection of feedback mechanisms -- such as AGN, star formation, and
radiative cooling in cluster cores -- to gas entropy, and the role
of this feedback in altering global properties of the ICM and 
impacting galaxy formation. I have also been studying a method for
quantifying the virialization state of clusters through the band
dependence of X-ray temperatures. For my thesis I assembled a sample
of 350 archival {\textit{Chandra}} observations for 276 clusters. The
results of this laborious task have been many and are detailed in the
research summary.

My expertise in X-ray astronomy, with {\textit{Chandra}}
specifically, ideally suits me to further study the evolution of
galaxy clusters and mass-observable relations as tools for
conducting tests of cosmological models and better understanding galaxy
formation/evolution. Adaptation of my skill sets to study clusters
in the radio, optical, and infrared is the next step in my
career and should come with a short learning curve thanks to my
existing, mature programming ability. I am a great asset for anyone
studying clusters both for my technical skills and to furthering
their research objectives. I feel the post-doctoral position at
Stanford is an excellent fit for me, and the groups research goals will
benefit from my addition.

Along with this letter are my curriculum vitae, a list of
publications, a summary of past and current research, and a brief
description of possible research directions. Letters of recommendation
from Megan Donahue, Mark Voit, and Jack Baldwin will arrive under
separate cover. Please do not hesitate to contact me if there is any
further information I can provide as you review my application.

Thank you for your consideration.

Sincerely,\\
\begin{minipage}{7.5in}
\epsfig{file=signature.ps,height=0.75in}
\end{minipage}
Kenneth W. Cavagnolo\\
Michigan State University

\end{document}
